(1) 点质量的横向引力偏折角有向形式为 \(\hat\alpha(\theta)=4GM/(c^2D_l\theta)\)。薄透镜几何给 \(\beta=\theta-(D_{ls}/D_s)\hat\alpha(\theta)\)，即 \(\beta=\theta-\theta_E^2/\theta\)，其中 \(\theta_E^2=4GM D_{ls}/(c^2D_lD_s)>0\)。这一角是长度比组合形成的角尺度，具体宇宙学中三距离应取适当角直径距离而不能直接以欧氏相减替代。

(2) 透镜方程化成 \(\theta^2-\beta\theta-\theta_E^2=0\)，解为 \(\theta_\pm=[\beta\pm\sqrt{\beta^2+4\theta_E^2}]/2\)。对 \(\beta>0\)，\(\theta_+>0,\theta_-<0\)，两个像位于透镜两侧。两像角距离 \(\theta_+-\theta_-=\sqrt{\beta^2+4\theta_E^2}\)，其乘积 \(\theta_+\theta_-=-\theta_E^2\)，可作为独立代数检验。

(3) 若源、透镜、观测者完美轴向对准 \(\beta=0\)，透镜方程给径向位置 \(|\theta|=\theta_E\)；绕对称轴旋转形成半径 \(\theta_E\) 的爱因斯坦环。在理想点源与严格轴对称点透镜模型下该环的切向放大率形式发散，但真实有限源尺寸及镜头分辨率使测量可见亮度有限；不能把数学奇性等同于无限总辐射能。

(4) 轴对称透镜的有符号放大率是 \(\mu=[(\beta/\theta)(d\beta/d\theta)]^{-1}\)。由 \(d\beta/d\theta=1+\theta_E^2/\theta^2\) 得 \(\mu_\pm=[1-(\theta_E/\theta_\pm)^4]^{-1}=\tfrac12[1\pm(\beta^2+2\theta_E^2)/(\beta\sqrt{\beta^2+4\theta_E^2})]\)。第二像为负奇偶性，观测总亮度放大取绝对值相加：令 \(u=\beta/\theta_E\)，则 \(\mu_{\rm tot}=(u^2+2)/[u\sqrt{u^2+4}]\ge1\)。当 \(u\gg1\)，\(\mu_{\rm tot}=1+2u^{-4}+O(u^{-6})\)，恢复远离透镜时放大趋1。
